\documentclass[12pt,a4paper]{article}
\usepackage[margin=2.5cm]{geometry}
\usepackage{amsmath,amssymb}
\usepackage{array}
\usepackage{booktabs}
\usepackage{multirow}
\usepackage{graphicx}
\usepackage{longtable}
\usepackage{hyperref}
\usepackage{physics}
\usepackage{chemformula}
\usepackage{upgreek}
\usepackage{enumitem}
\usepackage{float}
\usepackage{xcolor}
\usepackage{tikz}
\usepackage[T1]{fontenc}

\title{\textbf{Spectroscopic Term Symbols}}
\author{Manpreet Singh}
\date{}

\begin{document}
\maketitle
\tableofcontents
\newpage

%----------------------------------------------------------------------
\section{Quantum Numbers}
%----------------------------------------------------------------------

To completely describe an electron in an atom, four quantum numbers are needed: energy ($n$), angular momentum ($\ell$), magnetic moment ($m_\ell$), and spin ($m_s$).

The first quantum number describes the electron shell, or energy level, of an atom. The value of $n$ ranges from 1 to the shell containing the outermost electron.

\begin{table}[H]
\centering
\caption{Quantum Numbers}
\begin{tabular}{lll}
\toprule
\textbf{Number} & \textbf{Symbol} & \textbf{Possible Values} \\
\midrule
Principal Quantum Number & $n$ & $1, 2, 3, 4, \ldots$ \\
Angular Momentum Quantum Number & $\ell$ & $0, 1, 2, 3, \ldots, (n-1)$ \\
Magnetic Quantum Number & $m_\ell$ & $-\ell, \ldots, -1, 0, +1, \ldots, +\ell$ \\
Spin Quantum Number & $m_s$ & $+\tfrac{1}{2},\ -\tfrac{1}{2}$ \\
\bottomrule
\end{tabular}
\end{table}

\begin{table}[H]
\centering
\caption{Orbital Designations}
\begin{tabular}{ccccccc}
\toprule
$n$ & $\ell$ & $m_\ell$ & \textbf{No.\ of Orbitals} & \textbf{Orbital Name} & \textbf{No.\ of Electrons} \\
\midrule
1 & 0 & 0 & 1 & 1s & 2 \\
\midrule
\multirow{2}{*}{2} & 0 & 0 & 1 & 2s & 2 \\
 & 1 & $-1, 0, +1$ & 3 & 2p & 6 \\
\midrule
\multirow{3}{*}{3} & 0 & 0 & 1 & 3s & 2 \\
 & 1 & $-1, 0, +1$ & 3 & 3p & 6 \\
 & 2 & $-2,-1,0,+1,+2$ & 5 & 3d & 10 \\
\midrule
\multirow{4}{*}{4} & 0 & 0 & 1 & 4s & 2 \\
 & 1 & $-1,0,+1$ & 3 & 4p & 6 \\
 & 2 & $-2,-1,0,+1,+2$ & 5 & 4d & 10 \\
 & 3 & $-3,-2,-1,0,+1,+2,+3$ & 7 & 4f & 14 \\
\bottomrule
\end{tabular}
\end{table}

%----------------------------------------------------------------------
\section{Term Symbols}
%----------------------------------------------------------------------

\begin{quote}
``The symbol which indicates electronic configuration \& different energy states of atom is called term symbol.''
\end{quote}

Term symbols are represented as $^{2S+1}L_J$.

HN Russell and FA Saunders introduced a system to fully represent the position of electrons in an atom. According to this method, in a multi-electron system, many different arrangements of electrons are possible. A term symbol is given for the ground state arrangement which follows all the laws (Hund's, Pauli, Aufbau etc.). Term symbols are assigned using the R--S coupling scale or L--S coupling scale.

\subsection{The Russell--Saunders Coupling Scheme (R--S or L--S Coupling)}

This scheme applies to smaller atoms ($Z < 30$) having less interaction between spin and orbital angular momentum.

The term symbol for the ground state configuration is:
\[
^{2S+1}L_J
\]
where:
\begin{itemize}
  \item $2S+1$ is the spin multiplicity,
  \item $L$ is the orbital angular momentum,
  \item $J$ is the total angular momentum.
\end{itemize}

All terms are determined by calculating:
\begin{enumerate}
  \item \textbf{s--s coupling}, which gives total spin angular momentum ($S$)
  \item \textbf{l--l coupling}, which gives total orbital angular momentum ($L$)
  \item \textbf{L--S coupling}, which gives total angular momentum ($J$)
\end{enumerate}

\subsubsection{(1) s--s Coupling (S)}

The value of the resulting spin angular momentum is obtained from the unpaired electrons in the outer orbit of the atom.
\[
S = s_1 + s_2 + s_3 + \cdots
\]

\textbf{Example:} For $\mathrm{P}^3$ (three unpaired electrons):
\[
S = \tfrac{1}{2} + \tfrac{1}{2} + \tfrac{1}{2} = \tfrac{3}{2}
\]
Spin multiplicity: $2S+1 = 2\!\left(\tfrac{3}{2}\right)+1 = 4$

Paired electron spins cancel due to opposite spin.

\subsubsection{(2) l--l Coupling (L)}

Total orbital angular momentum is the sum of the orbital angular momentum $\ell$ of each electron:
\[
L = m_{\ell_1} + m_{\ell_2} + m_{\ell_3} + \cdots
\]

The value of $L$ is assigned using the following letter designations:

\begin{table}[H]
\centering
\begin{tabular}{ccccccccc}
\toprule
$L$ & 0 & 1 & 2 & 3 & 4 & 5 & 6 \\
\midrule
Term & S & P & D & F & G & H & I \\
\bottomrule
\end{tabular}
\end{table}

\textbf{Example:} For $\mathrm{d}^7$, with $m_\ell = +2,+1,0,-1,-2$:
\[
L = 2(+2) + 2(+1) + 1(0) + 1(-1) + 1(-2) = 4+2+0-1-2 = 3 \Rightarrow F\text{ term}
\]

\subsubsection{(3) L--S Coupling (J)}

The total orbital angular momentum ($L$) couples with the total spin angular momentum ($S$) to give total angular momentum ($J$):
\[
J = (L+S),\ (L+S-1),\ \ldots,\ |L-S|
\]

\textbf{Rule:}
\begin{itemize}
  \item For \textbf{less than half-filled} orbitals: use the \textbf{minimum} value of $J$.
  \item For \textbf{more than half-filled} orbitals: use the \textbf{maximum} value of $J$.
\end{itemize}

\textbf{Example:} For $\mathrm{P}^2$, $S=1$, $L=1$:
\[
J = (1+1), \ldots, (1-1) = 2, 1, 0
\]
Since $\mathrm{P}^2$ is less than half-filled, $J=0$.

%----------------------------------------------------------------------
\subsection{Worked Examples}
%----------------------------------------------------------------------

\subsubsection{Ground State Term for $\mathrm{P}^3$ Configuration}

Electron arrangement ($m_\ell = +1, 0, -1$, one electron each):
\[
L = 1(+1) + 1(0) + 1(-1) = 0 \Rightarrow S\text{ term}
\]
\[
S = \tfrac{1}{2}+\tfrac{1}{2}+\tfrac{1}{2} = \tfrac{3}{2}, \quad 2S+1 = 4
\]
\[
J = \left|0+\tfrac{3}{2}\right| \text{ to } \left|0-\tfrac{3}{2}\right| = \tfrac{3}{2}
\]
\[
\boxed{\text{Ground state term}: {}^4S_{3/2}}
\]

\subsubsection{Ground State Term for $\mathrm{d}^2$ Configuration}

Electron arrangement (two electrons at $m_\ell = +2, +1$):
\[
L = 1(+2)+1(+1) = 3 \Rightarrow F\text{ term}
\]
\[
S = \tfrac{1}{2}+\tfrac{1}{2} = 1, \quad 2S+1 = 3
\]
\[
J = (3+1) \text{ to } (3-1) = 4, 3, 2
\]
Less than half-filled $\Rightarrow$ minimum $J = 2$.
\[
\boxed{\text{Ground state term}: {}^3F_2}
\]

\subsubsection{Ground State Term for $^7\mathrm{N}$}

Configuration: $1s^2\,2s^2\,2p_x^1\,2p_y^1\,2p_z^1$
\[
L = 2(0)+1(+1)+1(0)+1(-1) = 0 \Rightarrow S\text{ term}
\]
\[
S = -\tfrac{1}{2}+\tfrac{1}{2}+\tfrac{1}{2}+\tfrac{1}{2} = \tfrac{3}{2}, \quad 2S+1 = 4
\]
\[
J = \left|0+\tfrac{3}{2}\right| \text{ to } \left|0-\tfrac{3}{2}\right| = \tfrac{3}{2}
\]
\[
\boxed{{}^4S_{3/2} \text{ (quartet S three-halves)}}
\]

\subsubsection{Ground State Term for $^{26}\mathrm{Fe}$}

Configuration: $[\mathrm{Ar}]\,3d^6\,4s^2$; considering 3d$^6$: ($m_\ell = +2,+1,0,-1,-2$, filled as $\uparrow\downarrow,\uparrow,\uparrow,\uparrow,\uparrow$)
\[
L = 2(+2)+1(+1)+1(0)+1(-1)+1(-2)+0 = 2 \Rightarrow D\text{ term}
\]
\[
S = 0+\tfrac{1}{2}+\tfrac{1}{2}+\tfrac{1}{2}+\tfrac{1}{2}+0 = 2, \quad 2S+1 = 5
\]
\[
J = |2+2| \text{ to } |2-2| = 4, 3, 2, 1, 0
\]
More than half-filled $\Rightarrow$ maximum $J = 4$.
\[
\boxed{{}^5D_4 \text{ (quintet D four)}}
\]

\subsubsection{Ground State Term for $\mathrm{d}^9$ Configuration}

3d$^9$: ($\uparrow\downarrow,\uparrow\downarrow,\uparrow\downarrow,\uparrow\downarrow,\uparrow$; $m_\ell = +2,+1,0,-1,-2$)
\[
L = 2(+2)+2(+1)+2(0)+2(-1)+1(-2) = 2 \Rightarrow D\text{ term}
\]
\[
S = \tfrac{1}{2}, \quad 2S+1 = 2
\]
\[
J = \left(2+\tfrac{1}{2}\right) \text{ to } \left(2-\tfrac{1}{2}\right) = \tfrac{5}{2},\ \tfrac{3}{2}
\]
More than half-filled $\Rightarrow$ maximum $J = \tfrac{5}{2}$.
\[
\boxed{{}^2D_{5/2}}
\]

\subsubsection{Finding $L$, $S$ and $J$ from a Given Term Symbol}

\textbf{For ${}^3F$:}
\[
2S+1 = 3 \Rightarrow S = 1 \text{ (2 unpaired electrons)}; \quad L = 3 \text{ (F term)}
\]
\[
J = (3+1),(3+0),(3-1) = 4, 3, 2
\]

\textbf{For ${}^1S$:}
\[
2S+1 = 1 \Rightarrow S = 0 \text{ (0 unpaired electrons)}; \quad L = 0 \text{ (S term)}
\]
\[
J = (0+0) = 0
\]

%----------------------------------------------------------------------
\section{Term Symbols for $\mathrm{p}^2$ Configuration (Carbon)}
%----------------------------------------------------------------------

The electronic configuration of C is $1s^2\,2s^2\,2p^2$. Since the $1s$ and $2s$ orbitals are fully filled, only the two $2p$ electrons contribute to the term symbols.

Three cases arise:

\subsection*{Figure 1: Electrons with Parallel Spins}
\begin{itemize}
  \item[(a)] $L = +1$, $S = 1$, $2S+1 = 3$
  \item[(b)] $L = 0$, $S = 1$, $2S+1 = 3$
  \item[(c)] $L = -1$, $S = 1$, $2S+1 = 3$
\end{itemize}
$L$ values $+1, 0, -1$ correspond to $L=1$ (P term). All arrangements give $2S+1=3$, hence: ${}^3P$

\subsection*{Figure 2 \& 3: Electrons with Opposite Spins}
\begin{itemize}
  \item Group (a): $L$ values $+2,+1,0,-1,-2$ $\Rightarrow$ $L=2$ (D term), $S=0$: ${}^1D$
  \item Group (b): $L = 0$ $\Rightarrow$ $L=0$ (S term), $S=0$: ${}^1S$
\end{itemize}

\textbf{Summary:} The terms for $\mathrm{p}^2$ are ${}^3P$, ${}^1D$, and ${}^1S$.

\subsection{Energy Ordering by Hund's Rule}

\begin{enumerate}
  \item The term with the \textbf{highest spin multiplicity} ($2S+1$) is most stable (lowest energy). ${}^3P$ has $2S+1 = 3$, so it is most stable.
  \item Among terms with the same $2S+1$, the one with \textbf{higher $L$} is more stable. ${}^1D$ ($L=2$) is more stable than ${}^1S$ ($L=0$).
\end{enumerate}

\[
{}^3P < {}^1D < {}^1S \quad \text{(increasing energy)}
\]

\subsection{J Energy Levels for $\mathrm{p}^2$}

\textbf{For ${}^3P$} ($L=1$, $S=1$):
\[
J = 2, 1, 0 \Rightarrow {}^3P_2,\ {}^3P_1,\ {}^3P_0
\]

\textbf{For ${}^1D$} ($L=2$, $S=0$):
\[
J = 2 \Rightarrow {}^1D_2
\]

\textbf{For ${}^1S$} ($L=0$, $S=0$):
\[
J = 0 \Rightarrow {}^1S_0
\]

Increasing order of energy:
\[
{}^3P_0 < {}^3P_1 < {}^3P_2 < {}^1D_2 < {}^1S_0
\]

Ground state Russell--Saunders term: ${}^3P_0$.

%----------------------------------------------------------------------
\section{Microstates}
%----------------------------------------------------------------------

In a given configuration, electrons can be arranged in many different ways that have slightly different energy. These different arrangements are called \textbf{microstates}.

\subsection{(1) Microstates from a Single Sub-shell}

\[
\text{Number of microstates} = \frac{n!}{r!\,(n-r)!}
\]

where $n = 2 \times (\text{number of orbitals in subshell})$ and $r = $ number of electrons.

\textbf{Example -- $\mathrm{p}^2$:} $r=2$, $n = 2\times3 = 6$
\[
\frac{6!}{2!\,4!} = 15
\]

\textbf{Example -- $\mathrm{d}^1$:} $r=1$, $n=2\times5=10$
\[
\frac{10!}{1!\,9!} = 10
\]

\begin{table}[H]
\centering
\caption{Number of Microstates for $p$, $d$, and $f$ Sub-shells}
\begin{tabular}{ccc}
\toprule
$P^n$ Config & $d^n$ Config & $f^n$ Config \\
\midrule
$P^1 = 6$  & $d^1 = 10$  & $f^1 = 14$ \\
$P^2 = 15$ & $d^2 = 45$  & $f^2 = 91$ \\
$P^3 = 20$ & $d^3 = 120$ & $f^3 = 364$ \\
$P^4 = 15$ & $d^4 = 210$ & $f^4 = 1001$ \\
$P^5 = 6$  & $d^5 = 252$ & $f^5 = 3003$ \\
$P^6 = 1$  & $d^6 = 210$ & $f^6 = 3432$ \\
           & $d^7 = 120$ & $f^7 = 3003$ \\
           & $d^8 = 45$  & $f^8 = 3003$ \\
           & $d^9 = 10$  & $f^9 = 2002$ \\
           & $d^{10} = 1$ & $f^{10} = 1001$ \\
           &             & $f^{11} = 364$ \\
           &             & $f^{12} = 91$ \\
           &             & $f^{13} = 14$ \\
           &             & $f^{14} = 1$ \\
\bottomrule
\end{tabular}
\end{table}

\subsection{(2) Microstates from Two Different Sub-shells}

Multiply the microstates of each sub-shell.

\textbf{Example -- $\mathrm{p}^1\mathrm{d}^1$:}
\[
\text{For } p^1\text{: } n=6,\ r=1 \Rightarrow 6; \quad \text{For } d^1\text{: } n=10,\ r=1 \Rightarrow 10
\]
\[
\text{Total microstates} = 6 \times 10 = 60
\]

\subsection{(3) Microstates from Term Symbol (without J)}

\[
\text{Microstates} = (2S+1)(2L+1)
\]

\textbf{Example -- ${}^2D$:} $2S+1=2$, $L=2$
\[
(2)(2\times2+1) = 2 \times 5 = 10
\]

\textbf{Microstates for $\mathrm{p}^2$ configuration:}
\begin{align*}
{}^3P: &\quad (3)(2\times1+1) = 3\times3 = 9 \\
{}^1D: &\quad (1)(2\times2+1) = 1\times5 = 5 \\
{}^1S: &\quad (1)(2\times0+1) = 1\times1 = 1 \\
\text{Total} &= 9+5+1 = 15 \checkmark
\end{align*}

\subsection{Pigeonhole Diagram for $\mathrm{p}^2$}

Total microstates of $\mathrm{p}^2 = 15$, grouped as:

\begin{table}[H]
\centering
\small
\begin{tabular}{c|ccccc|ccc|ccc|ccc|c}
\toprule
$m_\ell$ & a & b & c & d & e & f & g & h & i & j & k & l & m & n & o \\
\midrule
$-1$ & & & $\downarrow$ & $\downarrow$ & $\uparrow\downarrow$ &  & $\uparrow$ & $\uparrow$ &  & $\uparrow$ & $\uparrow$ & & $\downarrow$ & $\downarrow$ & \\
$0$ & & $\downarrow$ &  & $\uparrow$ &  & $\uparrow$ &  & $\uparrow$ & $\uparrow$ &  & $\downarrow$ & $\downarrow$ &  & $\downarrow$ & $\uparrow\downarrow$ \\
$+1$ & $\uparrow\downarrow$ & $\uparrow$ & $\uparrow$ &  & & $\uparrow$ & $\uparrow$ &  & $\downarrow$ & $\downarrow$ &  & $\downarrow$ & $\downarrow$ &  & \\
\midrule
$M_L$ & $+2$ & $+1$ & $0$ & $-1$ & $-2$ & $+1$ & $0$ & $-1$ & $+1$ & $0$ & $-1$ & $+1$ & $0$ & $-1$ & $0$ \\
$M_S$ & $0$ & $0$ & $0$ & $0$ & $0$ & $+1$ & $+1$ & $+1$ & $0$ & $0$ & $0$ & $-1$ & $-1$ & $-1$ & $0$ \\
$M_J$ & $+2$ & $+1$ & 0 & $-1$ & $-2$ & $+2$ & $+1$ & $0$ & $+1$ & $0$ & $-1$ & $0$ & $-1$ & $-2$ & $0$ \\
\midrule
& & & (A) &  & & \multicolumn{8}{c}{(B)} & & (C) \\
\bottomrule
\end{tabular}
\end{table}

\begin{itemize}
  \item \textbf{Group A} ($M_L = +2,+1,0,-1,-2$): $L=2$ (D), $S=0$, $J=2$ $\Rightarrow$ five ${}^1D_2$
  \item \textbf{Group B} ($M_L = +1,0,-1$, $M_S=+1$): $L=1$ (P), $S=1$, $J=2,1,0$ $\Rightarrow$ ${}^3P_2$, ${}^3P_1$, ${}^3P_0$
  \item \textbf{Group C} ($M_L=0$): $L=0$ (S), $S=0$, $J=0$ $\Rightarrow$ ${}^1S_0$
\end{itemize}

\subsection{Microstates for $\mathrm{d}^2$ Configuration}

$n = 10$, $r = 2$:
\[
\frac{10!}{2!\,8!} = 45
\]

Term symbols for $\mathrm{d}^2$: ${}^3F$, ${}^3P$, ${}^1G$, ${}^1D$, ${}^1S$

\begin{align*}
{}^3F: &\quad 3(2\times3+1) = 3\times7 = 21 \\
{}^3P: &\quad 3(2\times1+1) = 3\times3 = 9 \\
{}^1G: &\quad 1(2\times4+1) = 1\times9 = 9 \\
{}^1D: &\quad 1(2\times2+1) = 1\times5 = 5 \\
{}^1S: &\quad 1(2\times0+1) = 1\times1 = 1 \\
\text{Total} &= 21+9+9+5+1 = 45 \checkmark
\end{align*}

Ground state term symbols for $\mathrm{d}^2$ by Hund's rule:
\[
{}^3F_2 < {}^3F_3 < {}^3F_4 < {}^3P_0 < {}^3P_1 < {}^3P_2 < {}^1G_4 < {}^1D_2 < {}^1S_0
\]

\subsection{(4) Microstates from Term Symbol with J}

\[
\text{Microstates} = (2J+1)
\]

\textbf{Example -- ${}^2P_4$:} $J=4$
\[
(2\times4+1) = 9
\]

%----------------------------------------------------------------------
\section{Electronic Spectra of Transition Metal Complexes}
%----------------------------------------------------------------------

When electrons are promoted from one energy level to another, spectra arise. These are high-energy transitions. Much lower energy vibrational and rotational transitions always occur simultaneously, resulting in considerable broadening of electronic absorption bands. Band widths are commonly $1000$--$3000\ \mathrm{cm}^{-1}$.

Not all theoretically possible electronic transitions are observed. Selection rules distinguish between `allowed' and `forbidden' transitions.

\subsection{Selection Rules}

\subsubsection{(1) Spin Selection Rule}

During a transition, the total spin of the molecule (atom or ion) must remain constant.
\[
\Delta S = 0 \Rightarrow \text{allowed}; \quad \Delta S \neq 0 \Rightarrow \text{forbidden}
\]

Allowed: singlet$\to$singlet, doublet$\to$doublet, triplet$\to$triplet.\\
Forbidden: singlet$\to$doublet, doublet$\to$triplet, etc.

\subsubsection{(2) Laporte Orbital Selection Rule}

\[
\Delta\ell = \pm1 \Rightarrow \text{allowed (sharp band)}; \quad \Delta\ell \neq \pm1 \Rightarrow \text{forbidden}
\]

Allowed: s$\to$p, p$\to$s, p$\to$d, d$\to$p, d$\to$f, f$\to$d.

In terms of symmetry: $g \to u$ and $u \to g$ are allowed; $g \to g$ and $u \to u$ are forbidden.

\begin{table}[H]
\centering
\begin{tabular}{ccc}
\toprule
Orbital & $\ell$ & Symmetry \\
\midrule
s & 0 & gerade \\
p & 1 & ungerade \\
d & 2 & gerade \\
f & 3 & ungerade \\
\bottomrule
\end{tabular}
\end{table}

\subsection{Intensity of Electronic Transitions}

\begin{table}[H]
\centering
\caption{Approximate Molar Absorption Constants ($\varepsilon$)}
\begin{tabular}{lc}
\toprule
\textbf{Type of Transition} & \textbf{Approx.\ $\varepsilon$ (L mol$^{-1}$ cm$^{-1}$)} \\
\midrule
Spin forbidden, Laporte forbidden & $10^{-2}$ to $1.0$ \\
Spin allowed, Laporte forbidden & $1$ to $10$ \\
Spin allowed, Laporte forbidden (p \& d orbital overlap) & $10$ to $10^2$ \\
Spin allowed, Laporte forbidden (intensity stealing) & $10^2$ to $10^3$ \\
Spin allowed, Laporte allowed & $10^4$ to $10^5$ \\
\bottomrule
\end{tabular}
\end{table}

\subsection{Splitting of $\mathrm{d}^n$ Terms}

Each free ion term is affected by ligands in a complex. The effect depends on the geometry of the complex.

\begin{itemize}
  \item s-orbitals: spherically symmetrical, unaffected.
  \item p-orbitals: all three affected equally, no splitting.
  \item d-orbitals: split into two states in octahedral/tetrahedral fields.
  \item f-orbitals: split into three states.
\end{itemize}

\subsection{Transformation of Spectroscopic Terms to Mulliken Symbols}

\begin{table}[H]
\centering
\begin{tabular}{ccc}
\toprule
\textbf{Spectroscopic Term} & \textbf{Octahedral Field} & \textbf{Tetrahedral Field} \\
\midrule
S & $A_{2g}$ & $A_1$ \\
P & $T_{1g}$ & $T_1$ \\
D & $E_g + T_{2g}$ & $E + T_2$ \\
F & $A_{2g} + T_{1g} + T_{2g}$ & $A_2 + T_1 + T_2$ \\
\bottomrule
\end{tabular}
\end{table}

%----------------------------------------------------------------------
\section{Electronic Spectra of Specific Systems}
%----------------------------------------------------------------------

\subsection{$\mathrm{d}^1$ System}

In the free gaseous metal ion, five $d$ orbitals are degenerate and no $d$--$d$ transitions occur.

In octahedral complexes of Ti(III) (d$^1$ system) such as $[\mathrm{TiCl}_6]^{3-}$ and $[\mathrm{Ti(H_2O)_6}]^{3+}$, the $d$ orbitals split and the one electron occupies the lower $t_{2g}$ level. The only possible transition is to the $e_g$ level.

Absorption spectrum of $[\mathrm{Ti(H_2O)_6}]^{2+}$ shows one band with peak at $20300\ \mathrm{cm}^{-1}$.

Peaks for various Ti(III) complexes:
\begin{itemize}
  \item $[\mathrm{TiCl_6}]^{3-}$: $13000\ \mathrm{cm}^{-1}$
  \item $[\mathrm{TiF_6}]^{3-}$: $18900\ \mathrm{cm}^{-1}$
  \item $[\mathrm{Ti(H_2O)_6}]^{2+}$: $20300\ \mathrm{cm}^{-1}$
  \item $[\mathrm{Ti(CN)_6}]^{3-}$: $22300\ \mathrm{cm}^{-1}$
\end{itemize}

\subsection{$\mathrm{d}^9$ System}

$\mathrm{Cu}^{2+}$ ($Z=29$): $[\mathrm{Ar}]\,3d^9\,4s^0$, term symbol ${}^2D$.

In $[\mathrm{Cu(H_2O)_6}]^{2+}$, the octahedral field splits $d$ orbitals into lower ${}^2t_{2g}$ ($d_{xy},d_{yz},d_{xz}$) and higher ${}^2e_g$ ($d_{x^2-y^2},d_{z^2}$).

According to CFT, only one transition ${}^2T_{2g} \to {}^2E_g$ should be observed (one band). However, the spectrum of $[\mathrm{Cu(H_2O)_6}]^{2+}$ shows two broad, asymmetric bands. This is explained by the \textbf{Jahn--Teller effect}.

According to Jahn--Teller theory, a non-linear molecule containing degenerate orbitals in its ground state is distorted, destroying the degeneracy. Thus, $[\mathrm{Cu(H_2O)_6}]^{2+}$ loses its symmetry and the energy levels split further, giving three transitions:
\begin{enumerate}
  \item ${}^2B_{1g} \to {}^2E_g$
  \item ${}^2A_{1g} \leftarrow {}^2E_g$
  \item ${}^2B_{2g} \leftarrow {}^2E_g$
\end{enumerate}

The ${}^2B_{2g} \leftarrow {}^2E_g$ transition is very low energy (near IR), so only two bands are observed in the visible region. These two bands appear close together as a broad peak with a shoulder. The result is that $[\mathrm{Cu(H_2O)_6}]^{2+}$ appears blue-green in colour.

\subsection{Correlation Diagrams}

\subsubsection{$\mathrm{d}^1$ Correlation Diagram}

Ground term: ${}^2D$. Ground state configuration: $t_{2g}^1 e_g^0$; excited state: $t_{2g}^0 e_g^1$.
Transition: ${}^2T_{2g} \to {}^2E_g$, with energy equal to $\Delta_o$.

Relative to the ${}^2D$ free ion:
\[
{}^2T_{2g} \text{ stabilised by } {-0.4\Delta_o}; \quad {}^2E_g \text{ destabilised by } {+0.6\Delta_o}
\]

\subsubsection{$\mathrm{d}^9$ Correlation Diagram}

Ground term also ${}^2D$. Ground state: $t_{2g}^6 e_g^3$ (hole in $e_g$). Excited state: $t_{2g}^5 e_g^4$ (hole in $t_{2g}$). The energy ordering is reversed relative to $d^1$:
\[
{}^2E_g \text{ stabilised by } {-0.4\Delta_o}; \quad {}^2T_{2g} \text{ destabilised by } {+0.6\Delta_o}
\]

\subsection{Orgel Diagram of $\mathrm{d}^1$ -- $\mathrm{d}^9$}

Both $d^1$ and $d^9$ have ground state term ${}^2D$. In $d^1$ the electron transitions; in $d^9$ the positron (hole) transitions. Absorption bands of both $d^1$ and $d^9$ occur at equal frequencies.

In a tetrahedral field, the energy level diagram for $d^1$ complexes is the inverse of that in an octahedral field and similar to $d^9$.

\subsection{$\mathrm{d}^2$ Correlation Diagram}

V(III) octahedral complexes have the $d^2$ configuration. Ground state term: ${}^3F$. Excited states: ${}^3P$, ${}^1G$, ${}^1D$, ${}^1S$.

By the spin selection rule, ${}^3F \to {}^3P$ is allowed; ${}^1G$, ${}^1D$, ${}^1S$ are forbidden.

Three absorption bands arise from:
\begin{enumerate}
  \item ${}^3T_{1g}(F) \to {}^3T_{2g}(F)$: $17200\ \mathrm{cm}^{-1}$
  \item ${}^3T_{1g}(F) \to {}^3T_{1g}(P)$: $\}$ $25600\ \mathrm{cm}^{-1}$
  \item ${}^3T_{1g}(F) \to {}^3A_{2g}(F)$: $\}$
\end{enumerate}

Because the energy difference between transitions (ii) and (iii) is very small, they overlap to give a single peak. Only two peaks appear in the spectrum of $[\mathrm{V(H_2O)_6}]^{3+}$.

\subsection{$\mathrm{d}^8$ Correlation Diagram}

Ni(II) octahedral complexes have the $d^8$ configuration. Ground state term: ${}^3F$. Excited states: ${}^3P$, ${}^1G$, ${}^1D$, ${}^1S$.

In octahedral field, the ${}^3F$ term splits into lower energy ${}^3A_{2g}$ and higher energy ${}^3T_{2g}$ and ${}^3T_{1g}(P)$ levels.

Three transitions in $[\mathrm{Ni(H_2O)_6}]^{2+}$:
\begin{enumerate}
  \item ${}^3A_{2g} \to {}^3T_{2g}(F)$: $8700\ \mathrm{cm}^{-1}$
  \item ${}^3A_{2g} \to {}^3T_{1g}(F)$: $14500\ \mathrm{cm}^{-1}$
  \item ${}^3A_{2g} \to {}^3T_{1g}(P)$: $25300\ \mathrm{cm}^{-1}$
\end{enumerate}

All three transitions are energetically sufficient, so the visible spectrum of $[\mathrm{Ni(H_2O)_6}]^{2+}$ contains three peaks.

\subsection{$\mathrm{d}^2$--$\mathrm{d}^8$ Orgel Diagram}

Both $d^2$ and $d^8$ have ground state term ${}^3F$ and excited states ${}^3P$, ${}^1D$, ${}^1S$.

The splitting of $d^2$ in octahedral field equals that of $d^8$ in tetrahedral field. The splitting of $d^2$ in tetrahedral field equals that of $d^8$ in octahedral field.

\subsection{Electronic Spectra of High-Spin $\mathrm{d}^5$ Ions}

For high-spin $d^5$ ions, all possible $d$--$d$ transitions are spin-forbidden, since any transition must involve a reversal of spin. As a result, bands in the spectra of high-spin Mn(III) and Fe(III) complexes are very weak and the compounds are nearly colourless.

The four quartets ${}^4G$, ${}^4F$, ${}^4D$, ${}^4P$ involve the reversal of only one spin; the other seven states are doublets and doubly spin-forbidden. Up to ten extremely weak absorption bands may be observed in an octahedral field.

\subsection{$\mathrm{d}^3$--$\mathrm{d}^7$ Orgel Diagram}

Cr(III) octahedral complexes have $d^3$; Co(III) octahedral complexes have $d^7$. Ground state term: ${}^4F$; excited states: ${}^4P$, ${}^2G$, ${}^2F$, ${}^2D$, ${}^2P$.

By spin multiplicity rule, ${}^4F \to {}^4P$ is allowed; ${}^2G$, ${}^2F$, ${}^2D$, ${}^2P$ are forbidden.

\textbf{$[\mathrm{Co(H_2O)_6}]^{2+}$} shows three absorption bands:
\begin{enumerate}
  \item ${}^4T_{1g}(F) \to {}^4T_{2g}(F)$: $8000\ \mathrm{cm}^{-1}$
  \item ${}^4T_{1g}(F) \to {}^4A_{2g}(F)$: $19600\ \mathrm{cm}^{-1}$
  \item ${}^4T_{1g}(F) \to {}^4T_{1g}(P)$: $21600\ \mathrm{cm}^{-1}$
\end{enumerate}

\textbf{Tetrahedral $[\mathrm{CoCl_4}]^{2-}$} shows intense blue colour. Possible transitions:
\begin{enumerate}
  \item ${}^4A_2(F) \to {}^4T_1(P)$: $15000\ \mathrm{cm}^{-1}$
  \item ${}^4A_2(F) \to {}^4T_1(F)$: $5800\ \mathrm{cm}^{-1}$
  \item ${}^4A_2(F) \to {}^4T_2(F)$: $3500\ \mathrm{cm}^{-1}$
\end{enumerate}

Dissolving $\mathrm{CoCl}_2$ in water produces a pale pink solution of $[\mathrm{Co(H_2O)_6}]^{2+}$, whereas in alcohol the tetrahedral $[\mathrm{CoCl_2(CH_3CH_2OH)_2}]$ forms with intense blue colour.

The tetrahedral complex $[\mathrm{CoCl_4}]^{2-}$ shows much more intense $d$--$d$ bands than the octahedral $[\mathrm{Co(H_2O)_6}]^{2+}$ because the tetrahedral complex has no centre of symmetry, helping to overcome the $g \to g$ Laporte selection rule.

\subsection{Orgel Diagram of $\mathrm{d}^4$--$\mathrm{d}^6$}

Ground state term for both $d^4$ and $d^6$ is ${}^5D$. Excited states: ${}^3h$, ${}^3G$, ${}^3F$, ${}^3D$, ${}^3P$, ${}^1I$, ${}^1D$, ${}^1S$.

Only the ${}^5D$ transition is spin-allowed. The ${}^5D$ term splits into $T_{2g}$ and $E_g$.
\begin{itemize}
  \item High-spin $d^6$: electron promoted $T_{2g} \to E_g$ (same as $d^1$).
  \item $d^4$: positron (hole) promoted $E_g \to T_{2g}$ (same as $d^9$).
\end{itemize}

In a tetrahedral field, the subscript $g$ is omitted from the Orgel diagram.

The combined Orgel diagram for $d^1$, $d^9$, $d^4$, $d^6$ systems shows analogous splitting patterns.

%----------------------------------------------------------------------
\section{Jahn--Teller Effect}
%----------------------------------------------------------------------

The Jahn--Teller theorem explains why certain six-coordinate complexes possess distorted octahedral geometry.

\begin{enumerate}
  \item An octahedral complex has a regular shape when both $t_{2g}$ and $e_g$ sets of $d$-orbitals are symmetrically occupied.
  \item When the $t_{2g}$ orbitals are asymmetrically occupied, \textbf{slight distortion} occurs. This happens when the central metal ion contains $1$, $2$, $4$, or $5$ electrons in $t_{2g}$.
  \item When the $e_g$ orbitals are asymmetrically filled, \textbf{strong distortion} occurs, potentially changing the octahedral shape to tetragonal or even square planar. High-spin $d^4$ and $d^9$, and low-spin $d^7$, $d^8$, $d^9$ configurations lead to strong distortion.
\end{enumerate}

\subsection{Explanation for $\mathrm{Cu}^{2+}$ ($\mathrm{d}^9$)}

Configuration: $t_{2g}^6 e_g^3$. Two possible arrangements:
\begin{enumerate}
  \item $t_{2g}^6\ (d_{z^2})^2\ (d_{x^2-y^2})^1$
  \item $t_{2g}^6\ (d_{z^2})^1\ (d_{x^2-y^2})^2$
\end{enumerate}

Due to unsymmetrical $e_g$ occupation, ligands along the $z$-axis move away from the nucleus, while ligands in the $x$-$y$ plane move closer. Thus, $\mathrm{Cu}^{2+}$ complexes have \textbf{four short and two long} metal--ligand bonds.

Examples:
\begin{itemize}
  \item Cupric chloride crystal: four $\mathrm{Cl}^-$ at $2.30$~\AA, two $\mathrm{Cl}^-$ at $2.95$~\AA.
  \item Cupric fluoride crystal: four $\mathrm{F}^-$ at $1.93$~\AA, two $\mathrm{F}^-$ at $2.27$~\AA.
\end{itemize}

In the Jahn--Teller splitting: $\Delta_o \gg \delta_1 > \delta_2$.

%----------------------------------------------------------------------
\section{Spectrochemical Series}
%----------------------------------------------------------------------

The relative strength of different ligands can be determined by examining the electronic spectra of their complexes. When a stronger ligand replaces a weaker one, the absorption band shifts to \textbf{shorter wavelength} (blue shift, higher energy). A weaker ligand causes a \textbf{red shift} (longer wavelength, lower energy).

The spectrochemical series, in order of increasing ligand strength, is:

\[
I^- < Br^- < Cl^- < F^- < \text{O-donors} < H_2O < NH_3 < \text{N-donors} < CN^- \approx CO
\]

%----------------------------------------------------------------------
\section{Assignment}
%----------------------------------------------------------------------

\subsection*{Short Questions}

\begin{enumerate}
  \item Define term symbol.
  \item Give the value of $L$ for $s$, $p$, $d$, $f$, and $g$.
  \item What is spin multiplicity?
  \item Give Pauli's Exclusion Principle and Hund's Rule.
  \item Define: Total degeneracy.
  \item Write the rules for the ground state term symbol.
  \item Find the ground state term symbol for $^7\mathrm{N}$ ($Z=7$), $^6\mathrm{C}$ ($Z=6$), $\mathrm{V}^{3+}$ ($Z=23$), $\mathrm{Mn}^{2+}$ ($Z=25$), $\mathrm{Ni}^{2+}$ ($Z=28$), $\mathrm{Cu}^{2+}$ ($Z=29$).
  \item Find $L$, $S$, spin multiplicity, number of unpaired electrons and $J$ values of ${}^3F$, ${}^1S$, ${}^5D$.
  \item Calculate microstates for $p^3$ and $d^2$.
  \item Explain: Bands of $d$--$d$ transition are weaker.
  \item Explain: Identification of octahedral and tetrahedral complexes through electronic spectra.
  \item According to the Laporte rule, transition metal complexes should be colourless, but this is not true. Explain.
  \item Complexes of $\mathrm{Zn}^{2+}$ are mostly colourless. Explain.
  \item Explain: The electronic absorption band of \textit{cis}-$\mathrm{MA_4B_2}$ is more intense than \textit{trans}-$\mathrm{MA_4B_2}$.
\end{enumerate}

\subsection*{Long Questions}

\begin{enumerate}
  \item Explain the R--S (L--S) coupling scheme to derive the ground state term symbol.
  \item What is L--S coupling? Explain with the example of the $p^2$ configuration of carbon.
  \item Derive the ground state term symbol for: (i) $^7\mathrm{N}$, (ii) $^{16}\mathrm{S}$, (iii) $^{22}\mathrm{Ti}$, (iv) $\mathrm{Cr}^{3+}$ ($Z=24$), (v) $\mathrm{Fe}^{3+}$ ($Z=26$), (vi) $\mathrm{Cu}^{2+}$ ($Z=29$), (vii) $\mathrm{Co}^{2+}$ ($Z=27$), (viii) $\mathrm{Ni}^{2+}$ ($Z=28$).
  \item Identify the ground state term (giving reason) for: (i) ${}^1S$, (ii) ${}^3F$, (iii) ${}^3P$, (iv) ${}^1G$, (v) ${}^1D$.
  \item Calculate $L$ for $p^2$ configuration and show the energy states.
  \item Draw the pigeonhole diagram for $p^2$ configuration and arrange terms according to energy level.
  \item Draw the pigeonhole diagram for $d^2$ configuration and arrange terms according to energy level.
  \item Discuss the Selection Rules for electronic transition.
  \item Explain the Jahn--Teller effect in $[\mathrm{Ti(H_2O)_6}]^{3+}$.
  \item Discuss the electronic spectra of $[\mathrm{Ti(H_2O)_6}]^{3+}$ complex ion.
  \item Explain the Jahn--Teller effect in $[\mathrm{Cu(H_2O)_6}]^{2+}$.
  \item Explain: Orgel energy diagrams of $d^1$ and $d^9$ are opposite to each other.
  \item Write a short note on hole formalism.
  \item In $[\mathrm{V(H_2O)_6}]^{3+}$, three transitions are possible but only two peaks are observed. Explain.
  \item In $[\mathrm{Ni(H_2O)_6}]^{2+}$, three peaks are observed. Explain.
  \item Draw and explain the combined Orgel diagram of $d^3$ and $d^7$.
\end{enumerate}

\end{document}
=======
\documentclass[12pt]{article}

\usepackage[a4paper,margin=1in]{geometry}
\usepackage{amsmath}
\usepackage{setspace}
\usepackage{enumitem}

\title{\textbf{Preparation of 0.1 M Acetate Buffer (pH 5.0)}}
\author{}
\date{}

\begin{document}
	
	\maketitle
	
	\section*{Aim}
	
	To prepare 100 mL of 0.1 M acetate buffer of pH 5.0 using 0.1 M sodium acetate and 0.1 M acetic acid solutions.
	
	\section*{Required Solutions}
	
	\begin{itemize}
		\item 0.1 M Acetic acid solution (\(\mathrm{CH_3COOH}\))
		\item 0.1 M Sodium acetate solution (\(\mathrm{CH_3COONa}\))
		\item Distilled water (\(\mathrm{H_2O}\))
	\end{itemize}
	
	\section*{Apparatus}
	
	\begin{itemize}
		\item 100 mL volumetric flask
		\item Pipettes or measuring cylinders
		\item Beaker
		\item Glass rod
		\item pH meter
	\end{itemize}
	
	\section*{Theory}
	
	An acetate buffer consists of a weak acid (acetic acid, \(\mathrm{CH_3COOH}\)) and its conjugate base (sodium acetate, \(\mathrm{CH_3COONa}\)). It resists changes in pH when small amounts of acid or base are added. The pH of the buffer is governed by the Henderson--Hasselbalch equation:
	
	\[
	\mathrm{pH}=\mathrm{p}K_a+\log\left(\frac{[\mathrm{CH_3COO^-}]}{[\mathrm{CH_3COOH}]}\right)
	\]
	
	For acetic acid,
	
	\[
	\mathrm{p}K_a=4.76
	\]
	
	For a buffer of pH 5.0,
	
	\[
	5.0=4.76+\log\left(\frac{[\mathrm{CH_3COO^-}]}{[\mathrm{CH_3COOH}]}\right)
	\]
	
	Therefore,
	
	\[
	\frac{[\mathrm{CH_3COO^-}]}{[\mathrm{CH_3COOH}]}
	=10^{(5.0-4.76)}
	=10^{0.24}
	\approx1.74
	\]
	
	Since both stock solutions have the same concentration (0.1 M), the concentration ratio is equal to the volume ratio.
	
	Thus,
	
	\[
	\frac{V_{\mathrm{CH_3COONa}}}
	{V_{\mathrm{CH_3COOH}}}
	=1.74
	\]
	
	For a total volume of 100 mL,
	
	\[
	V_{\mathrm{CH_3COONa}}
	=
	\frac{1.74}{1+1.74}\times100
	=
	63.5~\mathrm{mL}
	\]
	
	\[
	V_{\mathrm{CH_3COOH}}
	=
	\frac{1}{1+1.74}\times100
	=
	36.5~\mathrm{mL}
	\]
	
	Hence, mixing 63.5 mL of 0.1 M sodium acetate with 36.5 mL of 0.1 M acetic acid produces a 0.1 M acetate buffer of pH 5.0.
	
	\section*{Procedure}
	
	\begin{enumerate}
		\item Clean a 100 mL volumetric flask and rinse it with distilled water.
		\item Pipette 63.5 mL of 0.1 M sodium acetate (\(\mathrm{CH_3COONa}\)) into the volumetric flask.
		\item Add 36.5 mL of 0.1 M acetic acid (\(\mathrm{CH_3COOH}\)).
		\item Mix the solution thoroughly by gentle swirling.
		\item Measure the pH using a calibrated pH meter.
		\item If necessary, adjust the pH:
		\begin{itemize}
			\item Add a few drops of 0.1 M acetic acid to decrease the pH.
			\item Add a few drops of 0.1 M sodium acetate to increase the pH.
		\end{itemize}
		\item If the total volume is slightly less than 100 mL after pH adjustment, make up the volume to the mark with distilled water.
		\item Mix well and label the solution as \textbf{0.1 M Acetate Buffer (pH 5.0)}.
	\end{enumerate}
	
	\section*{Observation}
	
	The prepared solution was clear and colourless. The measured pH was approximately 5.0.
	
	\section*{Result}
	
	A 100 mL sample of 0.1 M acetate buffer of pH 5.0 was successfully prepared by mixing 63.5 mL of 0.1 M sodium acetate solution with 36.5 mL of 0.1 M acetic acid solution.
	
\end{document}
>>>>>>> af2b87bca9c802b644b54b694f54664ddb47cf29
